近年来，随着微电子技术的快速发展，芯片热流密度不断攀升，高效散热问题引起了国内外学者的广泛关注，具有十分重要的理论意义和应用价值。微通道流动沸腾利用相变潜热，被认为是解决高热流密度散热问题最有前景的技术之一[1-3]。然而，流动沸腾不稳定性会诱发局部干涸，导致临界热流密度显著降低，严重制约了其工程应用。

申请人前期研究证明，在微通道底面构筑亲疏水相间的纳米结构可以有效抑制流动不稳定性：在质量流速为200 kg/(m2·s)的条件下，压降波动幅值从35 kPa降低到9 kPa，临界热流密度提高了约42%。本项目首次提出“润湿性梯度调控汽化核心分布”的新思路，将填补国内外在该领域的研究空白，研究成果有望达到国际领先水平，为高性能芯片散热器的设计奠定坚实的理论基础。
